\section{Appendice, dimostrazioni}

\subsection{Dimostrazione del Test $\chi^2$ (Goodness-of-Fit)}

\textbf{PRINCIPIO LOGICO:} confrontiamo la \textbf{distribuzione
osservata} dei dati (ciò che sta accadendo ora nella rete) con
una \textbf{distribuzione attesa} (un modello teorico di ciò che
consideriamo ``traffico normale'').

L'idea di base è verificare l'\textbf{Ipotesi Nulla ($H_0$)},
la quale afferma che ``il sistema si trova in uno stato normale''.
Se la discrepanza tra ciò che osserviamo e ciò che ci aspettiamo
è troppo grande, RIFIUTIAMO questa ipotesi e SEGNALIAMO
UN'ANOMALIA.

Per quantificare la discrepanza, usiamo la statistica del
$\chi^2$:
\[
    \chi^2 = \sum_{i} \frac{(O_i - E_i)^2}{E_i}
\]
\begin{itemize}
    \item varianza osservata rispetto all'attesa
    \item vogliamo valori positivi, non importa se il valore
          osservato è sopra o sotto quello atteso
    \item NORMALIZZAZIONE
\end{itemize}

\subsubsection{Esempio Pratico: Analisi delle Porte TCP}

Relativo al conteggio dei tentativi di connessione sulle varie
porte TCP in un istante $T$.

\textbf{K categorie:}
\begin{itemize}
    \item Porta 22 (SSH)
    \item Porta 80 (HTTP)
    \item Porta 443 (HTTPS)
    \item Porta 53 (DNS)
    \item Porta 3306 (MySQL)
    \item Porta 8080 (Admin)
    \item Altre porte (Other)
\end{itemize}

Abbiamo 7 categorie $\rightarrow$ \textbf{6 gradi di libertà
($d = k - 1$)}.

\textbf{Fasi:}
\begin{enumerate}
    \item \textbf{Stima del modello ($E_i$):} per ogni categoria,
    stimiamo quanti pacchetti ci aspettiamo in condizioni normali.
    Spesso, trattandosi di conteggi/arrivi nel tempo, si modellano
    i dati come variabili aleatorie di \textbf{Poisson} che, per
    grandi numeri, tendono a distribuzioni \textbf{Gaussiane}.
    Quindi $E_i$ rappresenta la nostra media storica.
    \[
        P(X=k) = \frac{\lambda^k e^{-\lambda}}{k!}
    \]
    \[
        E[X] = \sum_{k=0}^{\infty} k \cdot P(k) = \lambda
    \]
    \[
        Var(X) = E[X^2] - (E[X])^2
               = (\lambda^2 + \lambda) - \lambda^2 = \lambda
    \]
    \begin{itemize}
        \item Media $= N \cdot E_i$
        \item Varianza $\sigma^2 = E_i$
        \item Dividere per $E_i$ significa normalizzare l'errore
              rispetto alla varianza naturale del traffico.
    \end{itemize}

    \item \textbf{Osservazione ($O_i$):} contiamo effettivamente
    quanti pacchetti arrivano per ogni porta nell'intervallo
    corrente.

    \item \textbf{Calcolo del $\chi^2$:} sommiamo le differenze
    quadratiche normalizzate per tutte le 7 categorie.

    \item \textbf{La decisione: soglia e valori critici.} Una
    volta ottenuto il valore numerico $\chi^2$, per decidere se
    è un'anomalia dobbiamo confrontarlo con una
    \textbf{soglia critica}.
\end{enumerate}

\textbf{Limitazioni:}
\begin{itemize}
    \item Test applicabile solo a dataset monodimensionali.
    \item Modellare traffico di rete complesso con distribuzioni
          semplici e note è una sfida notevole. Se i dati non
          sono realmente ``poissoniani'' il test fallisce.
\end{itemize}

La \textbf{distribuzione teorica del $\chi^2$} è una funzione
di densità di probabilità nota, che dipende dai gradi di
libertà $d$.

Scegliamo un \textbf{livello di significatività $\alpha$}
(ad es. 0.05 o 0.01) che rappresenta la nostra tolleranza di
falsi positivi (quante occorrenze legittime siamo disposti a
classificare erroneamente come anomalie).

Confrontiamo le \textbf{tabelle tabulate} della distribuzione
$\chi^2$ (in letteratura) incrociando i gradi di libertà e
$\alpha$ per trovare il valore critico ($\chi^2_{critico}$):
\begin{itemize}
    \item Se $\chi^2 > \chi^2_{critico}
          \rightarrow$ ANOMALIA (rifiuto $H_0$)
    \item Se $\chi^2 < \chi^2_{critico}
          \rightarrow$ TRAFFICO NORMALE (accetto $H_0$)
\end{itemize}

\subsection{Ellipsoid Envelope}

Immagina di osservare un sistema che produce dati con 2
caratteristiche: il ``traffico di rete in ingresso'' e il
``carico della CPU''. Se visualizziamo questi dati in un
grafico cartesiano, vedremo una nuvola di punti. Nella
maggior parte dei sistemi stabili, questi punti tendono a
raggrupparsi in una zona centrale densa (comportamento normale),
mentre le anomalie si trovano lontane, isolate ai margini.

L'obiettivo dell'\textbf{Ellipsoid Envelope} è disegnare un
confine chiuso attorno a questa nuvola centrale. Dato che
assumiamo che i dati seguano una distribuzione Gaussiana
multivariata, la forma geometrica che meglio racchiude la
maggior parte dei punti è un'ellisse (2D) o un ellissoide
(3D o più). Tutto ciò che cade dentro questa forma è
``normale'' (``inlier''), tutto ciò che sta fuori è
``anomalia'' (``outlier'').

Per disegnare questa ellisse, l'algoritmo deve capire due cose:
dove si trova il centro della nuvola, quanto è grande e in che
direzione è orientata. Il centro dell'ellisse è determinato
dalla \textbf{media ($\mu$)} dei dati. La \textbf{forma e
l'orientamento} sono determinati dalla \textbf{matrice di
covarianza ($\Sigma$)}. Questa matrice è fondamentale perché
definisce non solo quanto i dati sono dispersi (varianza), ma
anche come le diverse caratteristiche si influenzano a vicenda
(covarianza). Se le variabili sono correlate l'ellisse non sarà
dritta, ma ruotata e schiacciata.

$\Sigma$ è complessa, ma attraverso la
\textbf{decomposizione spettrale} ($\Sigma = \Phi \Lambda
\Phi^T$), l'algoritmo la scompone in istruzioni geometriche
semplici:
\begin{itemize}
    \item La matrice $\Phi$ contiene gli \textbf{autovettori},
          che indicano le direzioni degli assi principali
          dell'ellisse. Ci dicono di quanti gradi dobbiamo
          ruotare la forma affinché si allinei con la
          distribuzione dei dati.
    \item La matrice $\Lambda$ contiene gli \textbf{autovalori},
          che indicano la lunghezza di questi assi. Un autovalore
          grande indica che in quella direzione i dati sono molto
          sparsi, quindi l'ellisse deve essere allungata; un
          autovalore piccolo indica che i dati sono concentrati,
          quindi l'ellisse deve essere stretta.
\end{itemize}

Vediamo $M$ come una matrice di trasformazione: immaginiamo
una sfera perfetta già standardizzata (media $= 0$ e
varianza $= 1$). Se applichiamo la trasformazione a questa
sfera, la ``stiriamo'' e ``ruotiamo'' fino a farla diventare
l'ellisse che si adatta ai nostri dati reali:
\[
    \Sigma = M \cdot M^T
\]

Per rilevare anomalie, l'algoritmo fa proprio l'inverso.
Prende un nuovo punto osservato e calcola la
\textbf{distanza di Mahalanobis} dal centro. Questa non è
una semplice distanza euclidea, ma una distanza ``pesata''
che utilizza l'inverso della covarianza ($\Sigma^{-1}$).
In termini geometrici, calcolare questa distanza equivale
ad applicare la trasformazione inversa ($M^{-1}$) al punto
osservato per vedere dove cadrebbe se l'ellisse fosse
riportata ad una sfera standard. Se il punto cade troppo
lontano dal centro (oltre una soglia definita dalla statistica
$\chi^2$ o dal rapporto di contaminazione), viene etichettato
come ANOMALIA.

\textbf{Limitazioni:} per calcolare correttamente $\mu$ e
$\Sigma$ in presenza di dati sporchi, l'algoritmo utilizza
stimatori robusti come \textbf{MCD (Minimum Covariance
Determinant)}. Questo metodo cerca il sottoinsieme di dati
più denso e compatto per calcolare l'ellisse, ignorando i
punti più distanti che potrebbero distorcere la forma.
La limitazione principale è che funziona bene solo su dati
``normali'' che hanno effettivamente una forma unimodale
(una singola bolla gaussiana). Se i dati hanno forme strane
o multiple, l'ellisse genera molti errori, poiché non si
adatta bene ai risultati.

\subsection{Tre Metodi/Livelli di Accesso al Modello}

\subsubsection{``Full Model'' (White-Box)}

Ha accesso completo ai pesi e all'architettura della rete
neurale $\rightarrow$ calcola facilmente il gradiente esatto
della funzione di costo $J$ rispetto all'input $x$.

\subsubsection{Metodo dei Logits (Prima di Softmax)}

È una variante del caso White-Box, ma suggerisce dove guardare
per ottenere un gradiente migliore.

\begin{itemize}
    \item \textbf{Logits ($z_i$):} sono i punteggi grezzi che
          escono dall'ultimo strato della rete neurale prima di
          essere convertiti in probabilità percentuali.
    \item \textbf{Softmax($z_i$):}
          \[
              \sigma(z)_i =
              \frac{e^{z_i}}{\sum_j e^{z_j}}
          \]
          Questa formula mostra come i logits vengono
          trasformati in probabilità.
\end{itemize}

Usare i logits è meglio della probabilità finale perché:
\begin{itemize}
    \item \textbf{Evitano la saturazione:} la funzione Softmax
          schiaccia i valori verso 0 o 1. Quando un valore è
          molto vicino a 1, la curva è piatta e il gradiente è
          quasi zero. I logits non hanno questo limite,
          fornendo un gradiente più vivo.
    \item \textbf{Preservano le distanze:} logits differenti
          possono produrre la stessa probabilità se differiscono
          per una costante (es. $z_i + C$), informazione che la
          Softmax cancella. I logits mantengono la ``superficie
          di separazione grezza'', dando più informazioni
          all'attaccante.
\end{itemize}

\subsubsection{Metodo dell'Approssimazione
            (Differenze Finite)}

Serve quando non possiamo calcolare il gradiente
(es. Black-box), ma possiamo solo interrogare il modello
e vedere il risultato (probabilità o loss).

Stima il gradiente ``sondando'' il modello con piccole
variazioni dell'input. Si utilizza il
\textbf{metodo delle differenze finite centrate}:
invece di calcolare la derivata analitica, si calcola la
pendenza della retta secante tra due punti vicinissimi.

\textbf{Formula A (basata sulla Loss):}
\[
    g_i \approx
    \frac{
        -\log p_y(x + \delta e_i)
        + \log p_y(x - \delta e_i)
    }{2\delta}
\]
con loss di cross-entropy $J = -\log(p)$.

\textbf{Cosa significa:} per capire quanto la feature $i$-esima
influenza l'errore, si prova a:
\begin{itemize}
    \item Aggiungere una piccola quantità $\delta$ a quella
          feature ($x + \delta e_i$)
    \item Sottrarre la stessa quantità ($x - \delta e_i$)
    \item Misurare la differenza nelle loss ($-\log p_y$) tra
          questi due punti
    \item Dividere per la distanza totale $2\delta$
\end{itemize}

\textbf{Risultato:} si ottiene un'approssimazione della
derivata parziale $\frac{\partial J}{\partial x_i}$ senza
conoscere la rete.

\textbf{Formula B (basata sulla Probabilità):}
\[
    \frac{\partial p_y}{\partial x_i}
    \approx
    \frac{
        p_y(x + \delta e_i) - p_y(x - \delta e_i)
    }{2\delta}
\]
con $J = -\log(p)$: se prob. $\uparrow$ $\rightarrow$ errore
$\downarrow$; se prob. $\downarrow$ $\rightarrow$ errore
$\uparrow$.

La derivata della probabilità indica fiducia $\rightarrow$
invertiamo il segno per ottenere l'errore.

\textbf{Formula di collegamento (Chain Rule):}
\[
    \frac{\partial J}{\partial x_i}
    = -\frac{1}{p_y(x)} \cdot
    \frac{\partial p_y}{\partial x_i}
\]

Un piccolo cambiamento nell'input consegue un'esplosione
nell'errore: se $p_y$ è molto basso (modello incerto) allora
$\frac{1}{p_y}$ è molto alto e amplifica. Questa formula serve
a convertire il risultato della Formula B (variazione della
probabilità) nel gradiente della loss necessario per FGSM.
Ci dice che il gradiente dell'errore è
\textbf{inversamente proporzionale alla probabilità della
classe corretta}.

\subsection{sFlow e NetFlow: Stima Statistica del Traffico}

\subsubsection{sFlow}

\[
    Rate = \frac{N}{n}
\]
Osservo $n$ pacchetti di $N$ che passano su un'interfaccia.
Sia $c$ il numero di pacchetti della classe di interesse,
$p$ l'incognita (proporzione reale della classe), allora:
\[
    N_c = n \cdot p, \qquad \sigma_c^2 = n \cdot p(1-p)
\]

Per una variabile lineare $y = Kx$:
\[
    E[y] = K E[x] = K\bar{x}
\]
\[
    Var(y) = K^2 \, Var(x)
\]

Posso stimare la proporzione incognita:
\[
    \hat{p} = \frac{c}{n}
    \quad \rightarrow \quad
    \text{proporzione incognita: quanti della frazione di }
    N \text{ appartengono alla classe } c.
\]

\[
    E[\hat{p}]
    = E\!\left[\frac{c}{n}\right]
    = \frac{1}{n} E[c]
    = \frac{1}{n} \cdot n \cdot p = p
\]

$\hat{p}$ è uno stimatore ottimo non polarizzato (unbiased).

\[
    Var(\hat{p})
    = Var\!\left(\frac{c}{n}\right)
    = \frac{1}{n^2} Var(c)
    = \frac{1}{n^2} \cdot n \cdot p(1-p)
    = \frac{p(1-p)}{n}
    \approx \frac{\hat{p}(1-\hat{p})}{n-1}
    \approx \frac{\frac{c}{n}\!\left(1-\frac{c}{n}\right)}{n-1}
\]

La utilizzo sui dati per sapere quanti pacchetti appartengono
alla classe $c$:
\[
    E[\hat{p}] = p = \frac{c}{n}, \qquad
    Var(\hat{p}) = \frac{c\!\left(1-\frac{c}{n}\right)}{n(n-1)}
\]

\subsubsection{NetFlow: Stima di $N_c$ e $B_c$}

NetFlow ci dice $N_c$, $B_c$ osservando $c$ pacchetti (byte
di flusso, numero di pacchetti appartenenti al flusso $c$
tra quelli che passano $N$).

\[
    E[N_c]
    = E[\hat{p} \cdot N]
    = N \, E[\hat{p}]
    = N \cdot \frac{c}{n}
\]
Numero medio di pacchetti appartenenti al flusso.

\[
    Var(N_c)
    = Var(\hat{p} \cdot N)
    = N^2 \, Var(\hat{p})
    = N^2 \,
      \frac{\frac{c}{n}\!\left(1-\frac{c}{n}\right)}{n-1}
\]

Posso determinare l'errore di stima per gli intervalli
di confidenza:
\[
    \left[
        N_c - z\sqrt{Var(N_c)},\;
        N_c + z\sqrt{Var(N_c)}
    \right]
\]

Lo standard error (SE) dipende in forma inversa dal numero
di campioni secondo la distribuzione Gaussiana. Per un
livello di confidenza del 95\% si ha $z = 1.96$.

L'errore standard è dell'ordine di $1/\sqrt{c}$:
\[
    \epsilon\%
    = 100 \cdot \frac{se}{N_c}
    = 100 \cdot 1{,}96 \, \frac{\sqrt{Var(N_c)}}{N_c}
    = 1{,}96 \sqrt{\frac{1}{c}\!\left(1-\frac{c}{n}\right)}
    < 1{,}96 \sqrt{\frac{1}{c}}
\]

con $c \leq n$: tra tutti i pacchetti osservati, al più
fanno tutti parte della stessa coorte, quindi
$(1 - \frac{c}{n}) < 1$.

\paragraph{Stima dei Byte $B_c$}

La variabile $B_c$ che avrei contato se avessi osservato
$N_c$ pacchetti, dove $\beta = b_1 + b_2 + \dots + b_c$ è
il numero totale di byte dei pacchetti appartenenti alla
coorte (variabili aleatorie statisticamente indipendenti
$\rightarrow Cov(\dots, \dots) = 0$):

\[
    B_c
    = \beta \cdot \frac{N_c}{c}
    = \left(\frac{1}{c}\sum_{i=1}^{c} b_i\right) N_c
    = \bar{b} \cdot N_c
\]

dove $\bar{b}$ è la media della dimensione dei pacchetti
appartenenti alla coorte $c$.

Ricordando la covarianza:
\[
    Cov(x,y)
    = E[(x-\bar{x})(y-\bar{y})]
    = E[xy] - \bar{x}\bar{y}
\]

e la varianza del prodotto di variabili indipendenti:
\[
    Var(x \cdot y)
    = Var(x)\,Var(y) + Var(x)\,\bar{y}^2
      + \bar{x}^2\,Var(y)
\]

La varianza della dimensione dei pacchetti è:
\[
    Var(b)
    = \frac{1}{c-1} \sum_{i=1}^{c} (b_i - \bar{b})^2
    = \frac{1}{c-1}
      \left[
          \sum_{i=1}^{c} b_i^2 - c\bar{b}^2
      \right]
    = \frac{1}{c-1}
      \left[
          \sum_{i=1}^{c} b_i^2
          - \frac{1}{c}
            \left(\sum_{i=1}^{c} b_i\right)^2
      \right]
\]

Varianza della dimensione dei pacchetti espressa solo in
base alla dimensione dei campioni.
